\documentclass{article}
\begin{document}
\title{An explicit state restore can reuse a body number without a world rewind}
\maketitle

\paragraph{Scope.}
Gazebo Sim 9 is pinned at
\texttt{62f198e0478159549b748c4b4c66f0cbe2fcbf6f}.
The test exercises an authorized state supplier in a serial graph fixture.
It runs the native state-bearing world-control callback, pending state
application, non-map \texttt{SetState} overload, entity creation, and
removal processing. Serialized components, full system restarts, joint
hierarchies, physical motion, external access policy, and real transport
are not in the test. Empty component lists and graph storage are
fixture inputs. The parent study and ordinary allocator guard assumed
the declared regular allocation path; this study adds a distinct
selected state application path.

\paragraph{Same accepted controls, different allocation lifetime.}
The first native allocation of the watched vertex has number 11.
The authorized supplier sends a state record deleting that number,
then a state record restoring it with that explicit identifier.
Both service replies are positive in both arms.
\begin{itemize}
\item If both inputs arrive before the next application of buffered
state, the second is the pending state that is applied. The vertex
remains present. No new lifetime is created by that application.
\item If one application/removal cut completes between the inputs,
the first input really removes the original. The second then creates
a new vertex at the supplied number 11. The native graph path to create
it does not go through ordinary monotonic allocation, and no world
rewind occurs.
\end{itemize}
The compared application cuts are paused; simulated time remains
constant and world-incarnation authority does not change.
The final selected membership and numeric identity agree. Thus a
receiver checking the current world generation in addition to the
number can still misapply an old allocation-scoped receipt. This
does not imply that an authorized semantic restore may never carry
an old duty forward; that requires a separately qualified map.

\paragraph{A source distinction and its lifetime.}
Creation marks the vertex as newly created. The pin clears that flag
in the step lifecycle after the selected source cut. A late poll for
the flag plus current ID can therefore agree with the coalesced arm.
Capture the actual birth at the source cut, and qualify the receipt
with a non-rewinding allocation lineage within the world generation
(or a fresh non-reusable operation key bound to that lineage).
For each tracked vertex the fixture increments a separate counter at
the actual birth witness, before it is cleared. A new allocation's
receipt can be admitted without invalidating the unaffected entities'
scope. Source-side observation must cover all permitted ways to
replace that vertex; starting a counter from an arbitrary late poll
does not supply that coverage. This is an authored projection; the
stock state callback does not provide a certified external issuance
history.

The replay compares 248 schedules of up to five such pairs and 1,032
pairs of admitted state inputs. An unrelated, separately numbered
allocation still executes at each round. Neither the authorized
body-state difference nor the supplier's cut is inferred from client
reply timing. A factored direct checker given the same source events
ties. No latency, cryptography, physical settling or cross-plant
proof-reuse claim is made.
\end{document}
